\documentclass[10pt]{article}
\usepackage[margin=1in, paperwidth=8.5in, paperheight=11in]{geometry}
\usepackage{ifpdf,amsmath, amssymb, comment, color, graphicx, stmaryrd,setspace,enumitem,tikz, fancyhdr, wrapfig, textcomp, units, mathptmx, pdfpages, cancel, multicol}
\usepackage{hyperref}
\hypersetup{
    colorlinks=true,
    urlcolor=blue,
}

\setlength{\headheight}{14.5pt}
\newcommand{\Q}{\mathbb{Q}}
\newcommand{\R}{\mathbb{R}}
\newcommand{\Z}{\mathbb{Z}}
\newcommand{\vu}{\mathbf{u}}
\newcommand{\vv}{\mathbf{v}}
\newcommand{\vw}{\mathbf{w}}
\newcommand{\vi}{\mathbf{i}}
\newcommand{\vj}{\mathbf{j}}
\newcommand{\vk}{\mathbf{k}}
\newcommand{\vn}{\mathbf{n}}
\newcommand{\vr}{\mathbf{r}}
\newcommand{\va}{\mathbf{a}}
\newcommand{\vF}{\mathbf{F}}
\newcommand{\vL}{\mathbf{L}}
\newcommand{\vT}{\mathbf{T}}
\newcommand{\vN}{\mathbf{N}}
\newcommand{\proj}{\operatorname{proj}}
\newcommand{\orth}{\operatorname{orth}}
\newcommand\del\nabla
\newcommand\dotp[1][.5]{\,\mathbin{\vcenter{\hbox{\scalebox{#1}{$\bullet$}}}}\,}
\newcommand{\slice}{_\text{slice}}

% Solution text is in red. If you want the solutions to show, remove the \iffalse from the definition of the \red command.
%\newcommand{\red}[1]{ %\iffalse
%	\textcolor{red}{#1} }%\fi}

\newenvironment{red}{\color{red}}{\ignorespacesafterend}
\newenvironment{blue}{\color{blue}}{\ignorespacesafterend}
\newcommand{\green}[1]{\textcolor{green}{#1}}
\renewcommand{\section}[1]{\begin{center} \textbf{#1} \\\end{center}}
%
\hyphenpenalty=5000
\setlength{\parindent}{0in}
%\oddsidemargin=-.25in
\allowdisplaybreaks
\pagestyle{fancy}
\renewcommand{\headrulewidth}{0pt}
\lhead{MATH 202}
\rhead{Spring 2024}
%\lfoot{\copyright\ CLEAR Calculus 2010}
\cfoot{}

\begin{document}
%


%\onehalfspacing
\allowdisplaybreaks
%##################################################################
\section{PS\#13 - Taylor series - \red{Answer key} }

\everymath{\displaystyle}
\begin{enumerate}[leftmargin=0pt]

\item (\href{https://faculty.gvsu.edu/boelkinm/Home/ACS/sec-8a-3-geometric.html#act-8a-3-3}{AC 8a.3 Activity 8.3.4}) Consider the function $f(x) = \dfrac{1}{1-x}$. Let's find the degree $n$ Taylor polynomial approximation centered at $a=0$ for $f$.
\begin{enumerate}
    \item To begin finding $T_n(x)$, do the usual work of computing the various derivatives of $f$ and their respective values at $a = 0$. Record the first five derivatives of $f(x)$ in the table below.
    \begin{align*}
        f(x) &= (1-x)^{-1} \\
        f'(x) &= (-1)(1-x)^{-2}(-1) \\
        f''(x) &= {\color{red} (-1)(-2)(1-x)^{-3}(-1)^2} \\
        f'''(x) &= {\color{red} (-1)(-2)(-3)(1-x)^{-4}(-1)^3} \\
        f^{(4)}(x) &= {\color{red} (-1)(-2)(-3)(-4)(1-x)^{-5}(-1)^4} \\
        f^{(5)}(x) &= {\color{red} (-1)(-2)(-3)(-4)(-5)(1-x)^{-6}(-1)^5}
    \end{align*}

    \item Next, evaluate the derivatives you determined in (a) at $a=0$ and use these to find the values of the coefficients of the Taylor polynomial centered at $a=0$.

    \begin{align*}
        f(0) &= (1-0)^{-1} = 1 & c_0 &= f(0) = 1 \\
        f'(0) &= (-1)(1-0)^{-2}(-1) & c_1 &= \dfrac{f'(0)}{1!} = \dfrac{1}{1!} = 1 \\
        f''(0) &= {\color{red} (-1)(-2)(1-0)^{-3}(-1)^2 = 2} 
        & c_2 &= \dfrac{f''(0)}{2!} = {\color{red} \dfrac{2}{2!} = 1}\\
        f'''(0) &= {\color{red} (-1)(-2)(-3)(1-0)^{-4}(-1)^3 = 6} 
        & {\color{red} c_3} &{\color{red} = \dfrac{f'''(0)}{3!} =  \dfrac{6}{3!} = 1}\\
        f^{(4)}(0) &= {\color{red} (-1)(-2)(-3)(-4)(1-0)^{-5}(-1)^4 = 24} 
        & {\color{red} c_4} &{\color{red} = \dfrac{f^{(4)}(0)}{4!} =  \dfrac{24}{4!} = 1}\\
        f^{(5)}(0) &= {\color{red} (-1)(-2)(-3)(-4)(-5)(1-0)^{-6}(-1)^5 = 120 }
        & {\color{red} c_5} &{\color{red} = \dfrac{f^{(5)}(0)}{5!} =  \dfrac{120}{5!} = 1}
    \end{align*}
    
    \item What pattern do you observe in the value of $c_k$? State the degree 5 Taylor polynomial, $T_5(x)$, as well as the formula you expect for the general degree $n$ Taylor polynomial, $T_n(x)$.
    
    \begin{red}
    Looks like $c_k$ is always 1. This makes sense because the value of $f^{(k)}(0)$ is always $k!$ based on how the negatives and powers work out.

    So, $T_5(x) = 1 + x + x^2 + x^3 + x^4 + x^5$, and $T_n(x) = 1 + x + x^2 + \ldots + x^n$. 
    \end{red}
    
    \item By identifying the value of $r$, explain why the degree $n$ Taylor polynomial $T_n(x)$ can be thought of as a finite geometric series.
    
    \begin{red}
    The value of $r$ is the ratio of the next term to the previous:
    \[r = \dfrac{c_{n+1}x^{n+1}}{c_n{x^n}} = \dfrac{1\cdot x^{n+1}}{1\cdot x^n} = x.\]

    At any particular point $x$, this is indeed a fixed ratio -- it doesn't depend at all on $n$.
    \end{red}

    \item What is the Taylor series centered at $a=0$ for $f(x) = \dfrac{1}{1-x}$?

    \begin{red}
        \[
        \dfrac{1}{1-x} = 1 + x + x^2 + x^3 + \ldots = \sum_{k=0}^{\infty} x^k.
        \]
    \end{red}
\end{enumerate}

\item (\href{https://faculty.gvsu.edu/boelkinm/Home/ACS/sec-8a-4-taylor-series.html#exercise-367}{AC 8a.4 Exercise 5}) As we will see in more detail in the next section, we can use known Taylor series to obtain other Taylor series, and we preview that idea in this exercise.

\begin{enumerate}
    \item Calculate the first four derivatives of $\sin(x^2)$ and hence find the fourth order Taylor polynomial for $\sin(x^2)$ centered at $a=0$.
    
    \begin{red}
    \begin{align*}
        f(x) &= \sin(x^2)\\
        f'(x) &= 2x \cos(x^2) \\
        f''(x) &= 2x(-\sin(x^2)\cdot 2x) + 2 \cos(x^2)  \\
        &= -4x^2 \sin(x^2) + 2\cos(x^2) \\
        f'''(x) &= -8x\sin(x^2) - 4x^2\cos(x^2)\cdot 2x + 2(-\sin(x^2)\cdot 2x) \\
        &= -8x^3 \cos(x^2) - 12x\sin(x^2) \\
        f^{(4)}(x) &= -24x^2\cos(x^2) -8x^3\cdot(-\sin(x^2)\cdot 2x)
        - 12 \sin(x^2) - 12x(\cos(x^2)\cdot 2x) \\
        &= -48x^2\cos(x^2)+(16x^4 - 12)\sin(x^2)
    \end{align*}
    At $a=0$:
    \begin{align*}
        f(0) &= \sin(0) = 0 \\
        f'(0) &= \cos(0) \cdot 2\cdot 0 = 0 \\
        f''(0) &= -4\cdot 0^2 \sin(0) + 2 \cos(0) = 2  \\
        f'''(0) &= -8\cdot 0^3\cos(0) - 12\cdot 0 \sin(0) = 0 \\
        f^{(4)}(0) &= -48\cdot 0^2 \cos(0) + (16\cdot 0^4 - 12)\sin(0) = 0
    \end{align*}
    So since a bunch of these derivatives are zero, $T_4(x) = 0 + 0x + \dfrac{2}{2!}x^2 + 0x^3 + 0x^4 $ -- that is, $T_4(x)$ is just $x^2$. (Boring!)
    \end{red}
    \item Part (a) demonstrates the direct approach to finding Taylor polynomials and series. Next we utilize a known Taylor series to make the process simpler. Recall that the Taylor series centered at 0 for $f(x) = \sin(x)$ is
    \begin{equation}
    T(x) = \sum_{k=0}^{\infty} (-1)^{k} \frac{x^{2k+1}}{(2k+1)!}\text{.}
    \end{equation}
    \begin{enumerate}
        \item[i.] Substitute $x^2$ for $x$ in the Taylor series $T(x)$ given above. Write out the first several terms and compare to your work in part (a). Explain why the substitution in this problem should result in the Taylor series for $\sin(x^2)$ centered at 0.

        \begin{red}
        \begin{align*}
            \text{Since } \sin(x) = T(x) &= \sum_{k=0}^{\infty} (-1)^{k} \frac{x^{2k+1}}{(2k+1)!}\text{,} \\
            \text{it makes sense that } \sin({\color{blue} x^2}) = T({\color{blue} x^2}) &= \sum_{k=0}^{\infty} (-1)^{k} \frac{({\color{blue} x^2})^{2k+1}}{(2k+1)!} \\
            &= \sum_{k=0}^{\infty} (-1)^{k} \frac{x^{\color{blue} 4k+2}}{(2k+1)!}
        \end{align*}
        The first few terms are:
        \[T({\color{blue}x^2}) = \underbrace{\dfrac{x^2}{1!}}_{k=0} - \underbrace{\dfrac{x^6}{3!}}_{k=1} + \underbrace{\dfrac{x^{10}}{5!}}_{k=2} - \underbrace{\dfrac{x^{14}}{7!}}_{k=3} - \ldots\]
        Makes sense that we didn't see anything past $x^2$ when we looked at just the first four derivatives.
        \end{red}

        \item[ii.] For what interval of $x$-values should we expect the Taylor series for $\sin(x^2)$ to converge?

        \begin{red}
            Since we came up with this series by substituting $x^2$ into the Taylor series for $\sin(x)$, which converges for all real values of $x$, I think this series should converge for all real values of $x$ also.
        \end{red}
    \end{enumerate}
    \end{enumerate}

\item (Bonus task!) In this problem we're going to explore the Taylor series for $f(x) = \sqrt{x}$ centered at $x = 1$.
\begin{enumerate}
    \item Find the Taylor polynomials for $f(x) = \sqrt{x}$ up through order 5, centered at $x=1$.
    \begin{red}
        \begin{align*}
            f(x) &= x^{1/2} 
            & f(1) &= 1 \\
            f'(x) &= \dfrac{1}{2} x^{-1/2} 
            &f'(1) &= \dfrac{1}{2} \\
            f''(x) &= \dfrac{1}{2}\left( -\dfrac{1}{2}\right) x^{-3/2} 
            & f''(1) &= \dfrac{1}{2}\left( -\dfrac{1}{2}\right) \\
            f'''(x) &= \dfrac{1}{2}
            \left( -\dfrac{1}{2}\right)
            \left( -\dfrac{3}{2}\right) 
            x^{-5/2} 
            & f'''(1) &= \dfrac{1}{2}
            \left( -\dfrac{1}{2}\right)
            \left( -\dfrac{3}{2}\right)  \\
            f^{(4)}(x) &= \dfrac{1}{2}
            \left( -\dfrac{1}{2}\right)
            \left( -\dfrac{3}{2}\right)
            \left( -\dfrac{5}{2}\right)
            x^{-7/2}
            & f^{(4)}(1) &= \dfrac{1}{2}
            \left( -\dfrac{1}{2}\right)
            \left( -\dfrac{3}{2}\right)
            \left( -\dfrac{5}{2}\right) \\
            f^{(5)}(x) &= \dfrac{1}{2}
            \left( -\dfrac{1}{2}\right)
            \left( -\dfrac{3}{2}\right)
            \left( -\dfrac{5}{2}\right)
            \left( -\dfrac{7}{2}\right)
            x^{-9/2}
            & f^{(5)}(1) &= \dfrac{1}{2}
            \left( -\dfrac{1}{2}\right)
            \left( -\dfrac{3}{2}\right)
            \left( -\dfrac{5}{2}\right)
            \left( -\dfrac{7}{2}\right) \\
        \end{align*}
        Okay, so $T_5(x) = 1 
        + \dfrac12 (x-1) 
        - \dfrac{1}{2!\cdot 2^2}(x-1)^2
        + \dfrac{1\cdot 3}{3!\cdot 2^3}(x-1)^3 
        - \dfrac{1\cdot3 \cdot 5}{4!\cdot 2^4}(x-1)^4
        + \dfrac{1\cdot3 \cdot 5\cdot 7}{5!\cdot 2^5}(x-1)^5$.
    \end{red}
    \item Spot the pattern in the coefficients.

    \begin{red}
        In the numerator we have an alternator $(-1)^{k+1}$ and the product of odd numbers. (If you want to read about ``double factorials'' that may be fun!) In the denominator we have a factorial and a power of 2:
        \[c_k = (-1)^{k+1}\dfrac{1\cdot 3 \cdot \ldots \cdot (2k-3)}{k!\cdot 2^k}.\]
    \end{red}
    \item Use the pattern you found to do the ratio test.
    \begin{red}
    \begin{align*}
        |r_n(x)| &= \left| \dfrac{c_{n+1}(x-1)^{n+1}}{c_n (x-1)^n} \right| \\
        &= \left|
        \dfrac{1\cdot 3\cdot \ldots\cdot(2n-3)(2n-1) (x-1)^{n+1}}{(n+1)!\cdot 2^{n+1}} \cdot        
        \dfrac{n!\cdot 2^n}{1\cdot 3\cdot \ldots\cdot(2n-3) (x-1)^n}
        \right| \\
        &= \dfrac{(2n-1)}{(n+1)\cdot 2}\cdot |x-1| \\
        |r(x)| &= |x-1| \lim_{n\to \infty} \dfrac{2n-1}{2n+2} = |x-1|\cdot 1
    \end{align*}
    \end{red}
    
    \item What is the ``interval of convergence'' -- the $x$-values for which the Taylor series will converge? Why does your answer make sense based on the graph of the original function $f(x) = \sqrt{x}$?

    \begin{red}
        This series will converge whenever $|r(x)| = |x-1| < 1$ -- that is, for $0 < x < 2$. This makes sense because the graph of $\sqrt{x}$ stops at $x=0$.
    \end{red}
    \item Use $T_5(x)$ to find a good fraction approximation for $\sqrt{2} \approx 1.414\ldots$

    \begin{red}
        \begin{align*}
            T_5(2) &= 1 
        + \dfrac12 (2-1) 
        - \dfrac{1}{2!\cdot 2^2}(2-1)^2
        + \dfrac{1\cdot 3}{3!\cdot 2^3}(2-1)^3 
        - \dfrac{1\cdot3 \cdot 5}{4!\cdot 2^4}(2-1)^4
        + \dfrac{1\cdot3 \cdot 5\cdot 7}{5!\cdot 2^5}(2-1)^5 \\
        &= 1 
        + \dfrac12 
        - \dfrac{1}{2!\cdot 2^2}
        + \dfrac{1\cdot 3}{3!\cdot 2^3}
        - \dfrac{1\cdot3 \cdot 5}{4!\cdot 2^4}
        + \dfrac{1\cdot3 \cdot 5\cdot 7}{5!\cdot 2^5} \\
        &= \dfrac{365}{256} = 1.42578125.
        \end{align*}
        Not bad!
    \end{red}
\end{enumerate}

    

\end{enumerate}
\end{document}